\section{Geometric Elegance versus Coordinate Practice }
\label{cm:geom}

\subsection{Introduction}

There are two ways to write down classical mechanics, and every
physicist learns them in the wrong order. First comes the practical
apparatus: pick generalized coordinates $q^i(t)$, write a Lagrangian
$L(q,\dot q,t)$ or a Hamiltonian $H(q,p,t)$, grind through the
Euler--Lagrange or Hamilton equations, solve. This is what one
actually computes with, and it is indispensable --- but it is also,
in a precise sense, an arbitrary choice of language: nothing about
the physics singles out these particular coordinates $q^i$ over some
other set. Only later, if at all, does a student meet the second
formulation: configuration space as an abstract smooth manifold,
velocities and momenta as elements of its tangent and cotangent
bundles, the dynamics generated by a connection or a symplectic
form defined without reference to any coordinate system whatsoever.

The two formulations describe exactly the same physical content; the
tension between them is not a tension between two theories but
between two languages for one theory. And as with any well-chosen
change of language, some things become obvious that were obscure,
and some things become laborious that were once immediate. This
chapter sets the two side by side, with a single worked example run
in both, and asks what each buys --- and costs.

\subsection{The Coordinate Form, and What It Hides}

The familiar machinery needs only a brief reminder. Generalized
coordinates $q^1,\dots,q^n$ parametrize the configurations of a
system; the Lagrangian $L(q,\dot q,t) = T - V$ is a function of
coordinates, velocities, and time; the equations of motion are
\begin{equation}
  \frac{d}{dt}\frac{\partial L}{\partial \dot q^i}
  - \frac{\partial L}{\partial q^i} \;=\; 0,
  \qquad i=1,\dots,n.
  \label{el-coord}
\end{equation}
This is the form in which mechanics is actually solved: one grinds
out~\eqref{el-coord} in whatever coordinates the problem hands
you, and out comes a trajectory.

Consider the simplest possible case: a free particle of mass $m$
moving in a plane, described in plane polar coordinates $(r,\phi)$
rather than Cartesian ones. The kinetic energy is $T =
\tfrac12 m(\dot r^2 + r^2\dot\phi^2)$, and there is no potential, so
$L=T$. Equation~\eqref{el-coord} gives
\begin{equation}
  m\ddot r - m r\dot\phi^2 \;=\; 0,
  \qquad
  \frac{d}{dt}\bigl(m r^2\dot\phi\bigr) \;=\; 0.
  \label{polar-free}
\end{equation}
The second equation is unremarkable --- it just says angular
momentum about the origin is conserved, as it must be for a free
particle. The first is the interesting one: it contains a term
$-mr\dot\phi^2$ that looks, for all the world, like an outward
centrifugal force. There is, of course, no force of any kind acting
on this particle; it travels in a perfectly straight line. The term
is not physics. It is the price of insisting on writing Newton's
second law in a coordinate system whose coordinate lines are
curved --- exactly analogous to the centrifugal and Coriolis terms
that appear when the same free-particle problem is worked in a
rotating frame, except that here the "frame" is not rotating at all;
only the \emph{coordinate grid} is curved.

This is not a quirk of polar coordinates specifically. In any
curvilinear coordinate system, extra terms of exactly this kind are
forced to appear, because the bare second derivative $\ddot q^i$ is
not, by itself, a coordinate-independent (tensorial) object; only the
combination
\begin{equation}
  \ddot q^i + \Gamma^i_{jk}(q)\,\dot q^j\dot q^k
  \label{covariant-accel}
\end{equation}
transforms correctly under a change of coordinates, where $\Gamma^i_{jk}$ are
the Christoffel symbols of the coordinate system (built from the kinetic-energy
metric $ds^2 = \sum_i m_i\,dq^i dq^i$). The Euler--Lagrange machinery
automatically assembles the correctly covariant
combination~\eqref{covariant-accel}, but it does so by generating, coordinate
by coordinate, extra terms that read exactly like additional applied forces ---
centrifugal, Coriolis, Euler --- even when, as here, the true physics is
completely force-free.

None of this is a defect of the coordinate approach; it is the
coordinate approach doing exactly what it is for. This is how
problems actually get integrated, in closed form or numerically, and
it is the only form in which mechanics ever makes contact with a
laboratory: no experiment outputs "a point of an abstract manifold,"
it outputs a number, which is to say a coordinate, read off a ruler,
a protractor, or a clock. A great deal of the technical skill of
mechanics --- finding action-angle variables for an integrable
system, Delaunay or Poincar\'e elements in celestial mechanics,
normal coordinates near an equilibrium --- consists precisely in
finding the coordinates in which a given problem's equations of
motion become as simple as possible. Coordinate choice is not
bookkeeping; it is frequently the whole content of solving a hard
problem.

But this leaves an uncomfortable question hanging. The same physical
system, worked in different coordinates, can produce equations of
motion that look wildly different in complexity, that display or
conceal conserved quantities, that contain "force" terms present in
one description and absent in another. Precisely because the
coordinate form does not, by itself, distinguish physics from
artifact, it cannot by itself answer a question that matters a great
deal: how much of what is written on the page in
equation~\eqref{polar-free} is actually about the world, and how
much is only about the choice of $(r,\phi)$?

\subsection{The Geometric Reformulation}

\subsubsection{Configuration space as a manifold}

The geometric answer to that question begins by refusing to
privilege any coordinate system in the first place. Configuration
space $Q$ is taken to be an abstract smooth manifold; a specific
choice of generalized coordinates $q^i$ is only ever a local chart on
$Q$, one among infinitely many equally valid alternatives, and
physical content is, by definition, whatever does not depend on
which chart one happens to be using.

Velocities of the system live in the tangent bundle $TQ$; the Lagrangian is
intrinsically a function $L : TQ\times\mathbb R \to \mathbb R$, stated without
reference to any coordinates at all. The Euler--Lagrange
equations~\eqref{el-coord} themselves have a coordinate-free meaning: they are
the flow of a specific vector field on $TQ$, picked out by the
Poincar\'e--Cartan structure that $L$ induces on $TQ$, satisfying the
second-order condition that this flow really does generate velocities
consistent with positions. What one writes down in coordinates
as~\eqref{el-coord} is the local expression of this single intrinsic object,
valid in every chart, looking different in each.

\subsubsection{Phase space and the symplectic form}

Passing (via the Legendre transform, where it is well defined) to
the momenta $p_i = \partial L/\partial \dot q^i$ carries the
description to the cotangent bundle $T^*Q$, the phase space. This
space carries a canonical structure that exists with reference to no
coordinates whatsoever: the \emph{tautological} (or Liouville) 1-form
$\theta$, defined at a point $(q,p)\in T^*Q$ by $\theta(v) =
p\bigl(\pi_*v\bigr)$ for any tangent vector $v$ there, where $\pi:
T^*Q\to Q$ is the projection. Its exterior derivative $\omega =
d\theta$ is the canonical symplectic form on $T^*Q$: closed,
nondegenerate, and built entirely out of the bundle structure of
$T^*Q$ over $Q$, with no reference to a coordinate system anywhere in
its definition.

Hamilton's equations are then the local expression of a single
intrinsic equation,
\begin{equation}
  \iota_{X_H}\,\omega \;=\; dH,
  \label{hamiltonian-vf}
\end{equation}
which defines the Hamiltonian vector field $X_H$ for any smooth
function $H$ on phase space. In local coordinates $(q^i,p_i)$ in
which $\theta = p_i\,dq^i$ and $\omega$ takes its standard constant
form, \eqref{hamiltonian-vf} reproduces exactly the familiar
\begin{equation}
  \dot q^i \;=\; \frac{\partial H}{\partial p_i},
  \qquad
  \dot p_i \;=\; -\,\frac{\partial H}{\partial q^i}.
  \label{hamilton-coord}
\end{equation}
Coordinates in which $\omega$ takes this simple constant form are
called \emph{canonical}, or \emph{Darboux}, coordinates, after Gaston
Darboux, who proved in 1882 that such coordinates always exist, at
least locally, around every point of every symplectic
manifold. This is a striking fact with no analogue
in Riemannian geometry, where curvature is a genuine local invariant
obstructing the metric from being brought to a standard flat form
even at a single point's neighborhood. Symplectic manifolds have no
such local invariant: every symplectic manifold of a given dimension
looks, locally, exactly like every other one. This is precisely why
Hamilton's equations can always be brought to the same simple
form~\eqref{hamilton-coord}, no matter how complicated the
physical system --- the complexity of a hard mechanics problem lives
entirely in the choice of \emph{which} canonical coordinates to use,
never in whether they exist.

\subsubsection{Canonical transformations reinterpreted}

In the coordinate treatment, canonical transformations tend to arrive
as a slightly mysterious piece of technology: a table of generating
functions $F_1,F_2,F_3,F_4$, and a set of conditions to check that a
proposed change of variables $(q,p)\mapsto(Q,P)$ preserves the form
of Hamilton's equations. From the geometric point of view, a
canonical transformation is simply a diffeomorphism of phase space
that preserves $\omega$ --- a \emph{symplectomorphism}. The
generating-function apparatus becomes a practical tool for
constructing particular symplectomorphisms, rather than the
definition of the concept. Darboux's theorem, in this light, is
exactly the statement that a symplectomorphism from a neighborhood
of any point to the flat model $(\mathbb R^{2n}, \sum dq^i\wedge
dp_i)$ always exists: "finding canonical coordinates" and "finding a
symplectomorphism to the standard model" are the same task, phrased
in two vocabularies.

\subsubsection{Geodesics, connections, and the return of the Christoffel symbols}

The same reinterpretation applies to the free-particle example. Motion
generated by a Riemannian metric --- whether the plain kinetic-energy metric of
a free system, or the Jacobi metric that reformulates a potential as pure
geometry --- is intrinsically \emph{geodesic} motion: $\nabla_{\dot x}\dot x =
0$, where $\nabla$ is the Levi-Civita connection of the metric, an object
defined without reference to any coordinate chart. Written out in coordinates,
this intrinsic equation reads
\begin{equation}
  \ddot x^i + \Gamma^i_{jk}(x)\,\dot x^j \dot x^k \;=\; 0,
  \label{geodesic-coord}
\end{equation}
with exactly the Christoffel symbols that appeared, unbidden, as
"fictitious centrifugal force" in equation~\eqref{polar-free}.
That is not a coincidence or an analogy: a free particle moving in a
straight line, described in polar coordinates, \emph{is} geodesic
motion in the flat metric $ds^2 = dr^2 + r^2 d\phi^2$, and the
Christoffel symbols of that metric in polar coordinates are precisely
the coefficients that turned up as a spurious force term. "Fictitious
force" and "Christoffel symbol" are the same object seen from two
sides. One calls it a force only for as long as one has declined to
think geometrically about it.

\subsection{What Each Form Buys You}

The geometric form earns its keep by telling you, in advance and for
free, which features of a set of equations of motion are physical
and which are artifacts of a chart. It explains why canonical
transformations exist and behave the way they do (Darboux's
theorem). It makes the relationship between symmetry and conservation
essentially automatic: a continuous symmetry is the flow of a vector
field preserving both $\omega$ and $H$, and when that symmetry comes
from the action of a Lie group, the associated conserved quantity is
the group's \emph{momentum map} --- the modern, coordinate-free
statement of Noether's theorem.
It also exposes genuinely global, topological facts that no single
coordinate patch can see: the phase space of a rigid body, for
instance, is the cotangent bundle of the rotation group $SO(3)$,
which cannot be covered by a single nonsingular coordinate chart at
all --- the root cause of gimbal lock, which textbooks often present
as an engineering nuisance but which is, geometrically, an
unavoidable topological obstruction. A coordinate singularity, on
this view, is never a physical singularity: it is only ever the
failure of one particular chart, repaired by changing coordinates or,
better, by working with the coordinate-free structure directly, where
the question of a "singularity" never even arises.

The coordinate form earns its keep by being the only form in which a
problem is actually solved, in closed form or on a computer, and the
only form that connects directly to a measurement, which is always
reported as a number relative to some frame. Darboux's theorem
guarantees good coordinates exist near any point; it says nothing
whatsoever about how to find the ones adapted to a \emph{specific}
Hamiltonian's structure --- action-angle variables for an integrable
system, say, which is frequently the entire technical content of
actually solving a hard problem in mechanics. No amount of abstract
symplectic geometry does that work for you.

Seen from the perspective of the wider project of this book, the two
forms sit on opposite sides of its recurring divide: a coordinate
expression is what gets compared, number for number, against
measurement, and so belongs to the "real world" side of the ledger;
the manifold, the connection, the symplectic form are pieces of
"ideal world" mathematics, whose entire justification is that they
correctly organize and predict the coordinate-level facts, and
nothing more. Neither side is dispensable, and neither is more
fundamental than the other in any sense that would let one be
discarded.

\subsection{A Worked Contrast: The Spherical Pendulum}

A single example run in both languages makes the trade-off concrete.
A pendulum of length $\ell$ and bob mass $m$, free to swing in any
direction, has configuration space the sphere of radius $\ell$
centered on the pivot.

\emph{Coordinate form.} Using the polar angle $\theta$ from the
downward vertical and the azimuth $\phi$, the Lagrangian is
\begin{equation}
  L \;=\; \tfrac12 m\ell^2\bigl(\dot\theta^2 + \sin^2\theta\,
  \dot\phi^2\bigr) \;+\; mg\ell\cos\theta .
  \label{sph-pendulum-L}
\end{equation}
The coordinate $\phi$ does not appear in $L$ itself, so its conjugate
momentum
\begin{equation}
  p_\phi \;=\; m\ell^2\sin^2\theta\,\dot\phi
  \label{sph-pendulum-pphi}
\end{equation}
is conserved --- it is the angular momentum about the vertical axis
--- and the problem reduces to an effective one-dimensional motion in
$\theta$ alone, with effective potential $V_{\mathrm{eff}}(\theta) =
mg\ell\cos\theta + p_\phi^2/(2m\ell^2\sin^2\theta)$. This is the form
in which the problem is actually solved, e.g.\ by quadrature or
numerically.

\emph{Geometric form.} Configuration space is the sphere $S^2$
itself, with the circle group $SO(2)$ acting on it by rotation about
the vertical axis, an action that preserves the Lagrangian. The
conservation of $p_\phi$ is nothing but the momentum map for this
$SO(2)$ symmetry, an instance of Noether's theorem requiring no
mention of $\theta,\phi$ at all. And the coordinate singularities of
$(\theta,\phi)$ at the poles $\theta=0,\pi$ --- where $\phi$ becomes
undefined and equation~\eqref{sph-pendulum-pphi} degenerates ---
are transparently artifacts of this particular chart: $S^2$ has no
distinguished points and no singularity there whatsoever. A pendulum
hanging straight down, or balanced straight up, is a perfectly smooth
physical configuration; it is only the coordinate system built from
$(\theta,\phi)$ that breaks down at exactly those two points, for the
same reason that all spherical coordinate systems break down at their
poles.

\subsection{Closing Remarks}

The two languages are not competitors; they specialize in different
questions. The geometric form is the place to ask what is really
happening, structurally, and what is guaranteed to remain true no
matter which coordinates one happens to choose. The coordinate form
is the place --- indeed the only place --- to ask what number a
system actually produces. A physicist who works only in coordinates
is at permanent risk of mistaking an artifact of a chosen chart for a
physical effect, exactly as happened, for a good part of the
twentieth century, with the coordinate singularity of the
Schwarzschild solution at its horizon: a singularity of the
particular $(t,r,\theta,\phi)$ coordinates originally used to write
the metric down, mistaken for decades for a genuine physical
singularity of spacetime itself, and only cleared up by passing to
coordinates (Eddington--Finkelstein, Kruskal) in which the horizon is
manifestly perfectly smooth --- a general-relativistic descendant, on
a much larger stage, of exactly the confusion embodied in the
centrifugal term of equation~\eqref{polar-free}. A physicist who
works only in the geometric language, on the other hand, has an
elegant theory of everything and the means to compute nothing. The
working practice of mechanics has always required both, spoken
fluently, and switched between as the problem at hand demands.



\subsection{Hertz's forceless mechanics }

In his \emph{Principles of Mechanics}, Heinrich Hertz reformulated dynamics
without the concept of force at all. Observable bodies are treated as coupled to
unobservable ``hidden masses'' through rigid, holonomic constraints, and the
combined system moves along a geodesic of a metric built from the mass
distribution on an enlarged configuration space. Suitably arranged, Hertz's
constrained geodesic motion reproduces the trajectories that Newtonian mechanics
attributes to the action of a force. The two formulations are empirically
equivalent: no experiment performed on the observable bodies alone can decide
between ``a force is acting'' and ``a constraint is enforcing geodesic motion in
a larger space.''

This is a case of underdetermination of theory --- indeed of ontology --- by
observation, arising entirely within classical mechanics. It is normally quantum
mechanics that is credited with forcing physicists to confront multiple,
empirically indistinguishable accounts of what is ``really'' happening beneath
the observed statistics: Bohmian trajectories, Everettian branches, and
epistemic and objective-collapse readings of the wavefunction are all consistent
with the same experimental record. Hertz shows that the same structural
situation --- one set of observations, several inequivalent stories about the
underlying ontology --- already existed in nineteenth-century particle mechanics,
with no probability amplitude in sight.


\chapter{Hertz's Mechanics Without Force}
\label{hm}

By the 1890s, several physicists --- Kirchhoff and Mach among them
--- had grown uneasy with the status of \emph{force} as a
fundamental concept of mechanics. Force is defined operationally
through $F=ma$, but that same equation is routinely used to
\emph{introduce} force, which makes the law read uncomfortably like
a definition dressed up as an empirical statement. In his
posthumously published \emph{Die Prinzipien der Mechanik in neuem
Zusammenhange dargestellt} (1894), Heinrich Hertz proposed to remove
the discomfort by removing force altogether: he set out to build the
whole of mechanics from just three primitive notions --- space,
time, and mass --- with no independent concept of force at
all~\cite{hertz}.

\section{The principle of the straightest path}

Consider $N$ mass points with positions $\mathbf x_i\in\mathbb R^3$
and masses $m_i$, and equip the $3N$-dimensional configuration space
$Q$ with the natural mass-weighted (kinetic-energy) metric,
\begin{equation}
  ds^2 \;=\; \sum_{i=1}^N m_i\, d\mathbf x_i\cdot d\mathbf x_i,
  \qquad
  T \;=\; \frac12\left(\frac{ds}{dt}\right)^{\!2}.
  \label{hertz-metric}
\end{equation}
Hertz's masses are never free in the naive sense: they are always
tied together by \emph{rigid, workless connections} (idealized
mechanical linkages that transmit force but do no work), which
restrict the accessible configurations to some submanifold
$\mathcal C\subset Q$. His fundamental postulate --- the
\emph{principle of the straightest path}
--- is then simply:

\begin{quote}
\emph{A system free of "force" moves along the straightest possible
path compatible with its constraints, at constant speed.}
\end{quote}

"Straightest possible" is to be read literally, in the geometric
sense: the representative point traces a \emph{geodesic} of the
metric~\eqref{hertz-metric} restricted to $\mathcal C$, i.e.\ a
curve whose acceleration (in the induced metric) is purely normal to
$\mathcal C$ and is entirely accounted for by the constraint
reaction. Formally, if the constraints are written $g_a(\mathbf
x)=0$, $a=1,\dots,k$, the equations of motion are
\begin{equation}
  m_i\,\ddot{\mathbf x}_i \;=\; \sum_a \lambda_a\,
  \frac{\partial g_a}{\partial \mathbf x_i},
  \label{hertz-eom}
\end{equation}
with the multipliers $\lambda_a$ fixed by demanding $g_a(\mathbf
x(t))\equiv 0$. There is nothing on the right-hand side of
\eqref{hertz-eom} except constraint reactions: no applied-force
term appears anywhere, by construction. Constant speed
($|ds/dt|=\mathrm{const}$, i.e.\ constant kinetic energy) follows
automatically, since ideal constraints do no work.

\section{Hidden masses, or where "force" goes}

Of course, ordinary bodies \emph{do} appear to exert forces on one
another --- gravity, contact forces, and so on all look like genuine
interactions between visible masses that are not rigidly connected
to anything. Hertz's resolution was to posit \emph{hidden masses}:
invisible mass points, rigidly linked to the visible ones (and, in
his more speculative moments, identified with a mechanical
aether), whose only role is to curve the constraint submanifold
$\mathcal C$ in just the right way. If the full configuration space
splits as $Q = Q_{\mathrm{vis}}\times Q_{\mathrm{hid}}$ and the
\emph{complete} system --- visible plus hidden --- moves along a
geodesic of $\mathcal C$, then projecting that geodesic motion down
onto $Q_{\mathrm{vis}}$ alone generically produces a curve that is
\emph{not} geodesic in the visible sector by itself: extra terms
appear in its equation of motion that are, from the point of view of
an observer with access only to $Q_{\mathrm{vis}}$, indistinguishable
from the effect of an applied force~\cite{hertz}.

This is a genuinely striking claim: on Hertz's account, force is not
absent from experience, only absent from the fundamental ontology.
Every observed force law would, in principle, be explained by some
specific (if permanently unobservable) arrangement of hidden mass
and rigid linkage. The same basic move --- explain an apparent force
as the shadow, on a submanifold, of pure geometric (geodesic) motion
in a larger space --- reappears twice in twentieth-century physics
in far more successful forms: in Kaluza--Klein theory, where the
electromagnetic force on a charged particle is recovered as geodesic
motion in a five-dimensional space, and, far more importantly, in
general relativity, where gravity itself --- which Newton had to
treat as a force acting at a distance --- becomes geodesic motion of
a free particle through curved four-dimensional spacetime, with no
"force" required at all. Hertz did not anticipate curved spacetime;
his geometry is Euclidean and the curvature is manufactured entirely
by hidden extra coordinates. But the structural resemblance to
Einstein's later, physically successful geometrization of gravity is
close enough that Hertz's book is still read today chiefly for that
reason, rather than as a working theory in its own
right.

\section{A second, more conservative geometrization}

It is worth noting that ordinary, force-\emph{having} mechanics can
also be written as a geodesic principle, without any hidden masses
at all, via Jacobi's form of the principle of least action. For a
conservative system with potential $V(\mathbf x)$ and fixed total
energy $E$, the true trajectories are exactly the geodesics of the
\emph{Jacobi metric}
\begin{equation}
  ds_J^2 \;=\; 2\bigl(E-V(\mathbf x)\bigr)\, ds^2 ,
  \label{jacobi-metric}
\end{equation}
a conformal rescaling of the mass metric~\eqref{hertz-metric} by
the local kinetic energy. So there are, in the end, two quite
different ways to eliminate force from the description of a system
governed by a potential: enlarge the configuration space with hidden
coordinates and keep the metric flat (Hertz), or keep the visible
configuration space as it is and let the metric itself become curved
and energy-dependent (Jacobi). Both reproduce exactly the same
observable trajectories. Neither is more "true" than the other; they
are, in Hertz's own vocabulary, different \emph{images} of the same
phenomena. That a concept as apparently basic as force can be made
to evaporate --- twice over, in two structurally different ways ---
without changing a single observable prediction is the real content
of this section: it shows "force" to be a feature of a chosen
representation of mechanics, not a fact about the world that every
correct representation must contain.

Hertz's own programme did not survive contact with its critics.
Almost every contemporary reviewer, Boltzmann included, objected
that positing a fresh, permanently unobservable mechanism of hidden
masses for every distinct force law (gravitation, electromagnetism,
\ldots) bought elegance at the price of testability, and the
research programme was essentially abandoned within a generation. It
survives today less as working physics than as an early, and
unusually explicit, demonstration that mechanics has more than one
internally consistent axiomatic foundation --- a lesson later put to
far more productive use.
